% Irodov Part IV §4.3 original 166--181. User PDF pp.187--188
% printed pp.178--179, answer PDF pp.322--323 (printed pp.321--323).
\waveproblem{166}{178}
长 $L=120$cm 的弦两端固定并形成驻波。
在弦上位移\emph{振幅}等于
$3.5$mm 的相邻点之间相距
$d=15.0$cm。
求弦振幅最大值，并指出所属泛音。
\exsolution
驻波振幅空间分布
$A(x)=A_{\max}|\sin kx|$。
方程 $|\sin kx|=3.5/A_{\max}$
在每个半波长周期内通常有两组解，相邻解点间隔交替。
题设把它们描述为相邻点间距均为同一个 $d$，
因此两类间距相等，解的相位必须处于
$|\sin kx|=1/\sqrt2$。
于是
\[
\boxed{A_{\max}=3.5\sqrt2\,\mathrm{mm}
\simeq5.0\,\mathrm{mm}}.
\]
这种情况下相邻指定点的相位差为 $\pi/2$，
$d=\lambda/4$，故 $\lambda=4d=0.60$m。
两端固定有
$L=n\lambda/2$，所以
$n=2L/\lambda=4$，
即第四谐波，
\[
\boxed{\text{第三泛音（third overtone）}}.
\]
\excheck “第三泛音”对应第四谐波而非第三谐波；
且数值 $5.0$mm 与原版答案一致。

\waveproblem{167}{178}
两根最初完全相同的弦，
分别被拉长
$\eta_1=2.0\%$、
$\eta_2=4.0\%$。
认为张力与相对伸长成正比。
求两根弦的基频之比。
\exsolution
设原弦长 $L_0$、原线密度 $\mu_0$、
材料张力常数为 $K$。
伸长比例 $\eta$ 后长度
$L=L_0(1+\eta)$，
张力 $T=K\eta$，
整根弦的质量不变，
线密度 $\mu=\mu_0/(1+\eta)$。
两端固定弦基频
$f=(1/(2L))\sqrt{T/\mu}$，
所以
$f\propto\sqrt{\eta/(1+\eta)}$。
因此
\[
\boxed{\frac{f_2}{f_1}
=\sqrt{\frac{\eta_2(1+\eta_1)}
{\eta_1(1+\eta_2)}}}
=\sqrt{\frac{0.04(1.02)}{0.02(1.04)}}
\simeq1.4.
\]
\excheck 此题不只是 $f\propto\sqrt T$，
因为伸长还改变振动部分的弦长及线密度。

\waveproblem{168}{178}
一根均匀张紧金属丝被缩短 $35\%$，
同时张力增加 $70\%$。
若弦材料的线密度保持不变，
求基频改变的方向及倍数。
\exsolution
弦两端固定基频
$f=(1/(2L))\sqrt{T/\mu}$，
缩短后
$L'=0.65L$，
新张力
$T'=1.70T$。
若取同材均匀金属丝重新截短而线密度不变，
\[
\boxed{\frac{f'}f=\frac1{0.65}\sqrt{1.70}
\simeq2.0}.
\]
因此基频\emph{升高至约2倍}。
\excheck 题目中“shortened”按截短同种金属丝理解；
如果是将原丝整体机械压短而改变密度/线密度，则还需额外模型。

\waveproblem{169}{178}
用一端为振动膜、另一端为可滑动活塞的声学共振管，
发声频率 $\nu=2000$Hz。
测得活塞的\emph{相邻共振位置}间距
$d=8.5$cm，求空气中的声速。
\exsolution
在相同端部边界条件下，
管长每增加半个波长，
便出现相邻一个纵向驻波共振：
$d=\lambda/2$。
故
$\lambda=2d=0.17$m，
声速
\[
\boxed{v=\nu\lambda=2\nu d
=2(2000)(0.085)
=340\,\mathrm{m/s}}.
\]
\excheck 不必知道端部修正的具体数值，
因为相邻共振位置的差将其消去。

\waveproblem{170}{178}
长 $L=85$cm 的空气柱中声速
$v=340$m/s。
统计低于
$\nu_{\max}=1250$Hz 的本征振动数：
(a) 一端封闭、一端开放；
(b) 两端均开放。
开放端视为位移腹点。
\exsolution
(a) 一端固定为位移结点，
另一端为位移腹点，
仅允许奇数四分之一波长：
\[
\boxed{\nu_n^{(a)}
=\frac{(2n+1)v}{4L}
=(2n+1)\,100\,\mathrm{Hz}},
\quad n=0,1,\ldots.
\]
小于1250Hz的为
100、300、500、700、900、1100Hz，
共
$\boxed{6\text{个}}$。
(b) 两端都是位移腹点：
\[
\boxed{\nu_n^{(b)}
=\frac{(n+1)v}{2L}
=(n+1)\,200\,\mathrm{Hz}},
\quad n=0,1,\ldots.
\]
小于1250Hz的为
200、400、600、800、1000、1200Hz，
也是
$\boxed{6\text{个}}$。
\excheck 频率上限是\emph{低于}1250Hz，
不包含等于上限的模态。

\waveproblem{171}{178}
铜棒长度 $L=50$cm，中点被夹紧，
求频率在20—50kHz之间的\emph{纵向}固有振动模式数及各个频率。
\exsolution
一维纵波声速
$v=\sqrt{E/\rho}$。
取铜的杨氏模量约
$E\simeq1.28\times10^{11}\,\mathrm{Pa}$、
密度约
$\rho=8.9\times10^3\,\mathrm{kg/m^3}$，
约得 $v\simeq3.79\,\mathrm{km/s}$，下文按书末约 $3.8\,\mathrm{km/s}$ 计算。铜的有效纵波速度依物性与试样而变，原题未给精确常数。
杆中点是位移结点，
两自由端为应力零、位移腹点；
所以每半根杆只能容纳奇数四分之一波长，
本征频率
\[
\boxed{\nu_n=\frac{(2n+1)v}{2L}
\simeq3.8(2n+1)\,\mathrm{kHz}},
\quad n=0,1,\ldots.
\]
满足 $20\le\nu_n\le50$kHz 的奇数倍依次是
7、9、11、13，因而
\[
\boxed{4\text{个模式： }
26.6,\ 34.2,\ 41.8,\ 49.4\ \mathrm{kHz}}.
\]
\excheck 使用原版书末近似 $v\simeq3.8$km/s，
模态编号 $n=3,4,5,6$；
原题没有指定铜的物性精确温度；旧稿误把 $E=1.1\times10^{11}\,\mathrm{Pa}$、$\rho=8.9\times10^3\,\mathrm{kg/m^3}$ 代成 $3.8\,\mathrm{km/s}$（该组数据实际约 $3.52\,\mathrm{km/s}$），此次改用与本页频率估计一致的有效模量。

\waveproblem{172}{178}
质量为 $m$ 的两端固定弦作\emph{基频}驻波，
角频率 $\omega$，最大位移振幅为 $a_{\max}$。
求：(a) 最大总动能；
(b) 一个完整周期的平均动能。
\exsolution
取弦长 $L$，
基频位移
$\xi(x,t)=a_{\max}\sin(\pi x/L)\cos\omega t$。
线密度
$\mu=m/L$，
总动能
\[
K(t)=\int_0^L\frac12\mu
(\partial_t\xi)^2\,\mathrm dx
=\frac12\frac mL a_{\max}^2\omega^2
\sin^2\omega t\int_0^L\sin^2(\pi x/L)\,\mathrm dx.
\]
积分等于 $L/2$，故
$K(t)=\frac14m a_{\max}^2\omega^2\sin^2\omega t$。
由此
\[
\boxed{K_{\max}=\frac14m\omega^2a_{\max}^2},
\quad
\boxed{\langle K\rangle=\frac18m\omega^2a_{\max}^2}.
\]
\excheck 平均 $\langle\sin^2\rangle=1/2$，
因而平均动能为最大值的一半。

\waveproblem{173}{178}
均匀长杆横截面积为 $S$、材料密度为 $\rho$，
内部纵向驻波
$\xi=a\sin kx\cos\omega t$。
求相邻两个位移结点之间包含的全部振动机械能。
\exsolution
结点间距为
$\Delta x=\pi/k$。
在 $t=0$ 质点速度处处为零，
机械能全部为弹性势能：
$w_p=\tfrac12\rho v^2(\partial_x\xi)^2
=\tfrac12\rho\omega^2a^2\cos^2kx$，
因为 $v=\omega/k$。
积分：
\[
W=S\int_0^{\pi/k}\frac12
\rho\omega^2a^2\cos^2kx\,\mathrm dx
=\frac12S\rho\omega^2a^2\frac{\pi}{2k}.
\]
因此
\[
\boxed{W=\frac{\pi S\rho\omega^2a^2}{4k}}.
\]
\excheck 对固定的一段相邻结点之间的介质，
无耗散时总机械能保持不变，故可在速度全零的时刻求能。

\waveproblem{174}{178}
频率 $\nu_0=1000$Hz 的声源以
$u=0.17$m/s 沿垂直墙面的直线运动。
两个静止接收器的空间排列为
$R_1$—声源—$R_2$—墙，
声速 $v=340$m/s。
指出哪个接收器会听到拍声以及拍频。
\exsolution
按题设声源向墙一侧移动。
$R_1$ 位于声源后方，接收到直接\emph{远离}声源的频率
$\nu_-= \nu_0v/(v+u)$，
以及在墙上反射、由源\emph{接近}墙面时确定的频率
$\nu_+=\nu_0v/(v-u)$。
因为墙静止，反射频率不变；
所以 $R_1$ 接收两不同频率的波，会产生拍声。
$R_2$ 位于声源前方，其直达与反射波具有相同源趋近频率，不产生这种拍声。
拍频为
\[
\boxed{\nu_b=\nu_+-\nu_-=
\frac{2\nu_0vu}{v^2-u^2}
\simeq\frac{2\nu_0u}{v}
\simeq1.0\,\mathrm{Hz}}.
\]
\excheck 判断接收器的关键是直接声与反射声所对应的源径向运动方向，而不是离墙远近。

\waveproblem{175}{179}
静止观察者同时听到两只频率均为
$\nu_0=680$Hz 的音叉，
一只以速度 $u$ 接近、另一只以同速远离；
听到拍频
$\nu_b=2.0$Hz。
空气声速 $v=340$m/s，求 $u$。
\exsolution
两只运动声源对应接收频率
$\nu_\pm=\nu_0v/(v\mp u)$，
拍频
\[
\nu_b=\nu_+-\nu_-
=\frac{2\nu_0vu}{v^2-u^2}.
\]
化为关于 $u$ 的二次方程，
$\nu_bu^2+2\nu_0vu-\nu_bv^2=0$，
取正根：
\[
\boxed{u=\frac v{\nu_b}
\left(\sqrt{\nu_0^2+\nu_b^2}-\nu_0\right)}.
\]
由于 $\nu_b\ll\nu_0$，
\[
\boxed{u\simeq\frac{v\nu_b}{2\nu_0}
=\frac{340(2)}{2(680)}
\simeq0.50\,\mathrm{m/s}}.
\]
\excheck 用数值计算直接相减两个相近大数容易舍入，
实际可采用等价表达式
$u=v\nu_b/(\sqrt{\nu_0^2+\nu_b^2}+\nu_0)$。

\waveproblem{176}{179}
频率 $\nu_0=2000$Hz 的声源沿接收者与声源连线作简谐往复运动，振幅 $a=50$cm、角频率为 $\Omega$；
静止接收者接收频率的上下限之差
$\Delta\nu=200$Hz，
声速 $v=340$m/s。
求 $\Omega$。
\exsolution
声源沿轴瞬时速度的极大模为
$u_{\max}=a\Omega$。
最高与最低接收频率分别为
$\nu_{\rm hi}=\nu_0v/(v-u_{\max})$、
$\nu_{\rm lo}=\nu_0v/(v+u_{\max})$。
故
\[
\Delta\nu=
\frac{2\nu_0v u_{\max}}{v^2-u_{\max}^2}.
\]
解出
$u_{\max}=v[\sqrt{\nu_0^2+\Delta\nu^2}-\nu_0]/\Delta\nu$，
进而
\[
\boxed{\Omega=\frac v{a\Delta\nu}
\left(\sqrt{\nu_0^2+\Delta\nu^2}-\nu_0\right)}.
\]
取 $\Delta\nu\ll\nu_0$ 近似：
\[
\boxed{\Omega\simeq\frac{v\Delta\nu}
{2a\nu_0}\simeq34\,\mathrm{s^{-1}}}.
\]
\excheck 瞬时频移由声源\emph{速度}决定，
不由位移振幅本身直接决定。

\waveproblem{177}{179}
频率 $\nu_0=1700$Hz 的声源与静止接收者最初位于同一点，
$t=0$ 起声源以
$w=10.0\,\mathrm{m/s^2}$ 的恒加速度远离接收者。
声速 $v=340$m/s。
求从开始运动算起 $t=10.0$s 时接收者\emph{真正听到}的频率。
\exsolution
接收时刻 $t$ 的声波在较早的发声时刻
$t_e$ 发出。
当时声源与接收者相距
$x_e=\tfrac12wt_e^2$，
声波传播历时 $x_e/v$，因此
\[
t=t_e+\frac{wt_e^2}{2v}.
\]
取物理解
$t_e=(v/w)[\sqrt{1+2wt/v}-1]$。
发声时刻的声源速度
$u_e=wt_e=v[\sqrt{1+2wt/v}-1]$。
接收者所听频率为
$\nu=\nu_0v/(v+u_e)$，故
\[
\boxed{\nu(t)=
\frac{\nu_0}{\sqrt{1+2wt/v}}}.
\]
代 $t=10$s，
\[
\boxed{\nu\simeq
\frac{1700}{\sqrt{1+200/340}}
\simeq1.35\,\mathrm{kHz}}.
\]
\excheck 若错误地使用\emph{接收时刻}的声源速度 $wt$，
而不考虑声波传播的迟延，会得到不同答案；
本题必须用发声时刻的速度。

\waveproblem{178}{179}
频率 $\nu_0=1.8$kHz 的点声源沿直线匀速运动，
其运动轨迹到静止观察者的最短距离为
$\ell=250$m，
声源速度 $u=\eta v$，
$\eta=0.80$，$v$ 为声速。
求：(a) 声源\emph{到达最近点的那个瞬间}观察者接收的频率；
(b) 观察者\emph{收到原频率} $\nu_0$ 时，声源与观察者的实际距离。
\exsolution
(a) 设声源最近点时刻 $t=0$，观察者位于 $(0,\ell)$；
此时收到的声波由较早的声源位置
$x_e=-ut_r$ 发出。
传播距离
$r_e=\sqrt{\ell^2+u^2t_r^2}=vt_r$，
解得
$t_r=\ell/[v\sqrt{1-\eta^2}]$。
传播方向上声源的趋近速度投影为
$u^2/v$，
因此接收频率
\[
\boxed{\nu_a=\frac{\nu_0v}{v-u^2/v}
=\frac{\nu_0}{1-\eta^2}
=\frac{1.8}{0.36}\,\mathrm{kHz}
=5.0\,\mathrm{kHz}}.
\]
(b) 要接收到原频率，必须对应\emph{声源恰经过最近点时}发出的声波；
该波抵达观察者需要 $\ell/v$ 时间。
其间声源又沿轨迹移动
$u\ell/v=\eta\ell$，
故接收瞬间的实际距离
\[
\boxed{r_b=\ell\sqrt{1+\eta^2}
=250\sqrt{1.64}\simeq0.32\,\mathrm{km}}.
\]
\excheck 本题(a)需考虑\emph{迟延发声位置}，
不可把当前瞬间径向速度零直接用于多普勒公式；
(b)源当前距离与发射时距离是不同的物理量。

\waveproblem{179}{179}
静止声源发出单色声波，
一面墙以速度
$u=33\,\mathrm{cm/s}$ 朝声源移动。
介质声速
$v=330\,\mathrm{m/s}$。
求反射声的波长相对入射波如何改变，百分比多少。
\exsolution
墙作为运动接收者先听到频率
$\nu_{\rm wall}=\nu_0(v+u)/v$；
反射时墙作为朝声源移动的运动发射者，
其反射波在介质中的频率
$\nu_{\rm ref}=\nu_{\rm wall}v/(v-u)
=\nu_0(v+u)/(v-u)$。
故波长比
\[
\frac{\lambda_{\rm ref}}{\lambda_0}
=\frac{\nu_0}{\nu_{\rm ref}}
=\frac{v-u}{v+u}.
\]
相对减少百分率
\[
\boxed{100\frac{\lambda_0-\lambda_{\rm ref}}
{\lambda_0}=100\frac{2u}{v+u}}
=\boxed{0.20\%}.
\]
\excheck $u=33$cm/s $=0.33$m/s，
所以 $2u/(v+u)\simeq0.0020$，
对应百分数\emph{0.20\%}；
本版书末似印2.0\%，有十倍出入，已记录。

\waveproblem{180}{179}
频率 $\nu_0=1700$Hz 的静止声源与静止接收者位于垂直于一面墙的直线上；
墙以速度
$u=6.0\,\mathrm{cm/s}$ 远离声源。
声速 $v=340$m/s。
求接收者听到的直达声与反射声之间的拍频。
\exsolution
墙作为向外远离的移动接收者，
接收频率
$\nu_{\rm wall}=\nu_0(v-u)/v$。
反射后墙相当于远离静止接收者的移动声源，
反射频率
\[
\nu_{\rm ref}=\frac{\nu_{\rm wall}v}{v+u}
=\nu_0\frac{v-u}{v+u}.
\]
直达频率仍为 $\nu_0$，
所以拍频
\[
\boxed{\nu_b=\nu_0-\nu_{\rm ref}
=\frac{2\nu_0u}{v+u}}.
\]
代 $u=0.060$m/s，
\[
\boxed{\nu_b\simeq0.60\,\mathrm{Hz}}.
\]
\excheck 对运动反射墙要连续应用两次多普勒关系，不能只做一次声源或接收者变换。

\waveproblem{181}{179}
各向同性点声源发声，在距声源
$r_1=10$m、$r_2=20$m 的两点，
声强之比为
$\eta=4.5$。
求声波振幅按
$e^{-\gamma r}$ 衰减时的衰减系数 $\gamma$。
\exsolution
球面波因几何扩散，位移幅度
$A(r)\propto e^{-\gamma r}/r$。
声强与幅度平方成比例，
故
$I(r)\propto e^{-2\gamma r}/r^2$。
因此
\[
\frac{I(r_1)}{I(r_2)}
=\left(\frac{r_2}{r_1}\right)^2
e^{2\gamma(r_2-r_1)}
=\eta.
\]
解出
\[
\boxed{\gamma=\frac1{2(r_2-r_1)}
\ln\left(\eta\frac{r_1^2}{r_2^2}\right)}.
\]
数据：
$\eta r_1^2/r_2^2=4.5/4=1.125$，
所以
\[
\boxed{\gamma=\frac{\ln1.125}{20\,\mathrm m}
\simeq5.9\times10^{-3}\,\mathrm{m^{-1}}
\approx0.006\,\mathrm{m^{-1}}}.
\]
\excheck 若忽略球面扩散，直接用声强比拟合指数衰减，
会高估介质本身的衰减系数。
